\phantomsection
\section*{ЗАКЛЮЧЕНИЕ}
\addcontentsline{toc}{section}{ЗАКЛЮЧЕНИЕ}

В ходе выполнения курсового проекта была рассмотрена предметная область
онлайн-бронирования столиков в ресторане, проведён обзор существующих
программных решений и определены требования к разрабатываемому
веб-приложению RestaurantBook.

В проектной части были выделены роли пользователей, описаны
пользовательские истории, построена диаграмма вариантов использования,
спроектирована модель данных и подготовлены шаблоны ключевых страниц.
Особое внимание уделено функциям, которые отличают разрабатываемую
систему от типовых решений: проверке доступности столиков, генерации
PDF-подтверждений, системе уведомлений и экспорту отчётов в форматах
CSV и PDF.

Результатом работы является проект веб-приложения, который может быть
использован как основа для дальнейшей реализации и развёртывания в
ресторане. Дальнейшее развитие проекта может включать интеграцию с
платёжными сервисами, расширенную аналитику посещаемости и адаптацию
интерфейса под мобильные устройства.


Разработанное веб-приложение было развёрнуто на платформе PythonAnywhere и доступно по адресу \url{https://anastasiaantipova.pythonanywhere.com/}. Исходный код проекта размещён в репозитории GitHub и предоставлен преподавателю для проверки. Развёртывание приложения на хостинге позволило проверить работоспособность системы в условиях, приближенных к реальной эксплуатации.
